% !PRINTED 652
\weberpage{652}
%<<CONT
gesetzt ist, und indem man für $\alpha_i$ in (7) die Ausdrücke (3) substituiert:
\weq{9}{\lambda \omega_r = \overset{i}{\Sigma} \omega_i \overset{s}{\Sigma} a_{s,i} t_{s,r}.}

Nach der Definition der Norm (§ 172) ist also $N(\lambda)$ die
Determinante aus den Koeffizienten
\[
\overset{s}{\Sigma} a_{s,i} t_{s,r},
\]
und diese kann man nach dem Multiplikationssatz der Determinante
 in
\weq{10}{N(\lambda) = A T}
zerlegen, worin $A$ die Bedeutung (4) hat, und $T$ die Determinante ist:
\weq{11}{T = \Sigma \pm t_{1,1} t_{2,2} \ldots t_{n,n}.}

Hierin ist $T$ eine homogene Funktion $n$ ten Grades der $t_i$,
deren Koeffizienten ganze Funktionen in $Z$ sind. $A$ ist selbst
eine ganze Funktion in $z$. Man beweist nun wie in Bd. II, § 164,
daß $T$ eine Einheit ist. Wäre das nämlich nicht der Fall, so
hätte $T$ irgend einen Linearfaktor $z - c$, und man könnte
nach einem elementaren Determinantensatz holomorphe ganze
Funktionale $y_i$ in $z$, die nicht alle durch $z - c$ teilbar sind, so
bestimmen, daß
\[
u_s = y_1 t_{s1} + \cdots + y_n t_{sn}
\]
durch $z - c$ teilbar wird (man kann den $y_i$ sogar konstanten
% [sic] printed konstanten Koeffizienten, the accusative plural object of geben without article requires konstante Koeffizienten (grammatical slip, mathematics unaffected)
Koeffizienten geben, da man sie auf ihre Reste nach $z - c$ reduzieren
 kann).

Setzt man nun
\[
\omega = \omega_1 y_1 + \omega_2 y_2 + \cdots + \omega_n y_n,
\]
so folgt aus (7)
\[
\lambda \omega = \alpha_1 u_1 + \alpha_2 u_2 + \cdots + \alpha_n u_n,
\]
und daraus folgt, weil die $\alpha_i$ durch $\lambda$, die $u_i$ durch $z - c$ teilbar
sind, daß $\omega$ durch $z - c$ teilbar ist, was der Definition der
Minimalbasis widerspricht. Demnach ergibt sich aus (10), daß $A$
die absolute Norm des Funktionals $\lambda$ ist:
\weq{12}{N_a(\lambda) = \Sigma \pm a_{11} a_{22} \ldots a_{nn}.}

Für ein Primfunktional $\pi$ können wir leicht eine Basis finden.
Es können nicht alle Elemente $\omega_1, \omega_2, \ldots, \omega_n$ einer Minimalbasis
durch $\pi$ teilbar sein, weil ja durch die Linearform
\weq{13}{\omega = x_1 \omega_1 + x_2 \omega_2 + \cdots + x_n \omega_n}
%CONT>>
